\chapter{Urea Breath Test}

\section*{What Is This Test?}
The urea breath test detects active infection with \textit{Helicobacter pylori}, a bacterium associated with peptic ulcer disease, chronic gastritis, and increased gastric-cancer risk. The patient ingests labeled urea. If \textit{H. pylori} urease is present in the stomach, it breaks down the urea and produces labeled carbon dioxide that is measured in exhaled breath.

\section*{Alternative Names}
\begin{itemize}
  \item H. pylori Breath Test
  \item Helicobacter pylori Urea Breath Test
  \item Carbon-13 Urea Breath Test
  \item Carbon-14 Urea Breath Test
\end{itemize}
\section*{Why It Is Ordered}
Initial noninvasive evaluation for active \textit{H. pylori} infection in appropriate patients; confirmation that eradication treatment was successful; and evaluation of recurrent symptoms when active infection is suspected. It detects active infection rather than simply past exposure.

\section*{How This Test Is Performed}
A baseline breath sample is collected, the patient drinks a labeled urea solution, and one or more additional breath samples are collected after the specified interval. The laboratory measures labeled carbon dioxide and reports the result as positive or negative according to the assay cutoff.

\section*{Preparation, Risks, and Limitations}
Proton-pump inhibitors, bismuth, and antibiotics can suppress \textit{H. pylori} and cause false-negative results. The ordering clinician will provide washout instructions; do not stop prescribed medicine without guidance. Fasting is commonly required. The test is noninvasive and generally safe, but a positive result should be interpreted with symptoms and the treatment history. Test-of-cure timing is important because testing too soon after treatment can be misleading.

\section*{References}
\begin{enumerate}
  \item MedlinePlus. Helicobacter pylori (H. pylori) Tests.
  \item American College of Gastroenterology. Clinical guideline for treatment of Helicobacter pylori infection.
\end{enumerate}

\clearpage
\section*{Understanding Results and Follow-up}
The laboratory report should identify the specimen, method, units, reference interval or decision limit, and whether the result is positive, negative, abnormal, or inconclusive. Reference intervals vary with age, sex, pregnancy, specimen type, assay, and laboratory; a result outside a range is not automatically a diagnosis, and a result within a range does not always exclude disease. Discuss unexpected results with the ordering clinician, who can decide whether repeat or confirmatory testing, treatment, or referral is appropriate. Do not change medicines or delay urgent care based on an isolated result.


\section*{Regulatory Status and Authorization Date}
This chapter describes a test category, not a specific commercial product. FDA, CDSCO, EU IVDR, and MHRA approval or registration dates therefore cannot be assigned without the assay manufacturer, product name, instrument, and intended use. For laboratory-developed tests, an individual agency approval date may not exist. See the Preface for the interpretation of regulatory status.

\clearpage
